%% toc: P3 35 35.1,35.2,35.3,35.4
\section{Comparison games and infinitary languages}
\label{sec:p03-w03-p03-038-comparison-games}

A comparison game tests whether two structures can be told apart by a
bounded amount of quantification. Spoiler selects elements, Duplicator
answers in the other structure, and Duplicator survives as long as the
selected pairs form a partial isomorphism. The length of the game measures
the quantificational resources: a fixed finite number of rounds matches
first-order quantifier rank, a game of length \(\omega\) matches the
infinitary language \(L_{\infty\omega}\), an ordinal clock matches
infinitary quantifier rank, and a game of length \(\omega_1\) can separate
structures that every \(L_{\infty\omega}\)-sentence confuses. This section
proves those correspondences and turns each into an inexpressibility
argument with explicit structures and explicit round bounds.

\emph{Inputs.} Quantifier rank and the finite game of
\cref{def:p03-quantifier-rank,def:p03-ef-game}; the infinitary language and
potential isomorphisms of \cref{def:p08-bvb501-potential-isomorphism} with
\cref{thm:p08-bvb501-potential-isomorphism}; the ordinal back-and-forth
approximants of \cref{def:p03-morley-approximants}; the countable dense
linear order theorems \cref{thm:p03-rn-dlo-extension,thm:p03-rn-dlo-uniqueness};
clubs and stationary sets in \cref{def:p06-club-stationary}; and, for the
source bridge, the presentation of Fra\"iss\'e's 1954 article in
\cref{sec:p03-fraisse-1954}.

\emph{Conclusions.} The finite-round theorem \cref{thm:p03-ef} is proved
for positions with parameters (\cref{thm:p03-w03-p03-038-pointed-ef}),
with the finiteness and definability lemmas it needs; the game criterion
and the connectivity application are proved with the explicit bound
\(m_k=(k+1)2^{k+2}\) (\cref{thm:p03-w03-p03-038-connectivity}); finite
back-and-forth systems are matched with the game, and Fra\"iss\'e's
extension condition is bridged to it; the closed game of length
\(\omega\), the clock games and the game of length \(\omega_1\) receive
their comparison lemmas; and a club and stationary-set analysis of one
family of dense orders of size \(\aleph_1\) shows the long game separating
\(L_{\infty\omega}\)-equivalent structures
(\cref{thm:p03-w03-p03-038-phi-games}).

\emph{Standing conventions.} The metatheory is ZFC. The Axiom of Choice
is used to assemble strategies from pointwise choices of answers and, in
\cref{sec:p03-w03-p03-038-clubs}, through the regularity of \(\omega_1\);
each use is named where it occurs. Structures have nonempty domains and
equality is logical. Signatures are relational. The finite-round results
of \cref{sec:p03-w03-p03-038-finite-games} require a finite signature;
the infinite games of \cref{sec:p03-w03-p03-038-long-games} do not.

\subsection{Quantifier rank and finite Ehrenfeucht--Fra\"iss\'e games}
\label{sec:p03-w03-p03-038-finite-games}

The game of \cref{def:p03-ef-game} starts from an empty position. Its
induction requires positions in which some pairs have already been
selected, so the first convention adds them.

\begin{hbconvention}{def:p03-w03-p03-038-positions}{Pointed games}
Fix a finite relational signature \(\tau\), \(\tau\)-structures
\(\mathcal A,\mathcal B\), and \(n\in\N\). Fix distinct variables
\(x_1,x_2,\ldots\). The \emph{atomic list} \(\mathrm{At}_n\) consists of
the formulas \(x_i=x_j\) for \(1\le i,j\le n\) and
\(R(x_{i_1},\ldots,x_{i_r})\) for every \(r\)-ary \(R\in\tau\) and
\(1\le i_1,\ldots,i_r\le n\). Because \(\tau\) is finite and has no
function or constant symbols, \(\mathrm{At}_n\) is finite.

A \emph{position} of length \(n\) is a pair of tuples
\(\bar a=(a_1,\ldots,a_n)\in(\dom\mathcal A)^n\) and
\(\bar b=(b_1,\ldots,b_n)\in(\dom\mathcal B)^n\), read as the pairing
\(a_i\mapsto b_i\). Repetitions are allowed. The pairing is a
\emph{partial isomorphism} when, for every \(\alpha\in\mathrm{At}_n\),
\(\mathcal A\models\alpha[\bar a]\) iff \(\mathcal B\models\alpha[\bar b]\);
this is exactly the condition of \cref{def:p03-ef-game}, which respects
repetitions and preserves and reflects every relation symbol.

In the \emph{pointed game} \(G_k(\mathcal A,\bar a;\mathcal B,\bar b)\),
\(k\) rounds are played as in \cref{def:p03-ef-game}, each round
appending its pair to the position; Duplicator wins a play exactly when
the final position, of length \(n+k\), is a partial isomorphism. A
\emph{strategy} for Duplicator assigns an answer to every position reached
after a Spoiler move; it is \emph{winning} if every play following it is won
by Duplicator. The game of \cref{def:p03-ef-game} is
\(G_k(\mathcal A,\varnothing;\mathcal B,\varnothing)\), written
\(G_k(\mathcal A,\mathcal B)\).

Since a restriction of a partial isomorphism to fewer coordinates is
again a partial isomorphism, Duplicator wins a play iff every intermediate
position of the play is a partial isomorphism.

For a formula \(\varphi\) whose free variables are among
\(x_1,\ldots,x_n\), write \(\mathcal A\models\varphi[\bar a]\) for
satisfaction under an assignment sending \(x_i\) to \(a_i\). Write
\((\mathcal A,\bar a)\equiv_k(\mathcal B,\bar b)\) when the two pointed
structures agree on every such formula of quantifier rank at most \(k\).
For \(n=0\) this is agreement through rank \(k\) in the sense of
\cref{def:p03-quantifier-rank}.
\index{Ehrenfeucht--Fraisse game@Ehrenfeucht--Fra\"iss\'e game!pointed}
\end{hbconvention}

\begin{hbdefinition}{def:p03-w03-p03-038-records}{Rank records}
With \(\tau\) as in \cref{def:p03-w03-p03-038-positions}, define for every
\(\tau\)-structure \(\mathcal A\) and \(\bar a\in(\dom\mathcal A)^n\) the
\emph{rank-\(j\) record} \(\operatorname{rec}_j(\mathcal A,\bar a)\) by
recursion on \(j\):
\[
 \operatorname{rec}_0(\mathcal A,\bar a)\defeq
   \{\alpha\in\mathrm{At}_n:\mathcal A\models\alpha[\bar a]\},
\]
\[
 \operatorname{rec}_{j+1}(\mathcal A,\bar a)\defeq
   \bigl(\operatorname{rec}_0(\mathcal A,\bar a),\;
   \{\operatorname{rec}_j(\mathcal A,\bar ac):c\in\dom\mathcal A\}\bigr),
\]
where \(\bar ac\) is the tuple of length \(n+1\) extending \(\bar a\) by
\(c\). Let \(\mathrm{Rec}_{n,0}\) be the power set of \(\mathrm{At}_n\)
and
\(\mathrm{Rec}_{n,j+1}\defeq\mathcal P(\mathrm{At}_n)\times
\mathcal P(\mathrm{Rec}_{n+1,j})\). Then
\(\operatorname{rec}_j(\mathcal A,\bar a)\in\mathrm{Rec}_{n,j}\) for all
structures and tuples, by induction on \(j\).
\index{rank record}
\end{hbdefinition}

\begin{hblemma}{thm:p03-w03-p03-038-finite-records}{proved=annex:p03-w03-p03-038-finite-records}{Finitely many records}
For every \(n\) and \(j\), the set \(\mathrm{Rec}_{n,j}\) is finite. If
\(N_n=|\mathrm{At}_n|\), then \(|\mathrm{Rec}_{n,0}|=2^{N_n}\) and
\(|\mathrm{Rec}_{n,j+1}|=2^{N_n}\cdot2^{|\mathrm{Rec}_{n+1,j}|}\).
\end{hblemma}

\begin{hbdefinition}{def:p03-w03-p03-038-characteristic}{Characteristic formulas}
For \(t\in\mathrm{Rec}_{n,j}\) define a formula
\(\chi_t(x_1,\ldots,x_n)\) by recursion on \(j\). For
\(t=S\subseteq\mathrm{At}_n\),
\[
 \chi_S\defeq\bigwedge_{\alpha\in S}\alpha\;\wedge
             \bigwedge_{\alpha\in\mathrm{At}_n\setminus S}\neg\alpha .
\]
For \(t=(S,Q)\in\mathrm{Rec}_{n,j+1}\),
\[
 \chi_t\defeq\chi_S\;\wedge\;
   \bigwedge_{q\in Q}\exists x_{n+1}\,\chi_q(x_1,\ldots,x_{n+1})
   \;\wedge\;
   \forall x_{n+1}\bigvee_{q\in Q}\chi_q(x_1,\ldots,x_{n+1}).
\]
Finite conjunctions and disjunctions are iterated binary ones, an empty
conjunction is \(\neg\bottom\) and an empty disjunction is \(\bottom\).
All index sets are finite by \cref{thm:p03-w03-p03-038-finite-records},
so \(\chi_t\) is a first-order formula.
\end{hbdefinition}

\begin{hblemma}{thm:p03-w03-p03-038-characteristic}{proved=annex:p03-w03-p03-038-characteristic}{Characteristic formulas define their records}
For \(t\in\mathrm{Rec}_{n,j}\), the free variables of \(\chi_t\) are among
\(x_1,\ldots,x_n\) and \(\operatorname{qr}(\chi_t)\le j\). For every
\(\tau\)-structure \(\mathcal A\) and \(\bar a\in(\dom\mathcal A)^n\),
\[
 \mathcal A\models\chi_t[\bar a]
 \iff
 \operatorname{rec}_j(\mathcal A,\bar a)=t .
\]
\end{hblemma}

\begin{hbtheorem}{thm:p03-w03-p03-038-pointed-ef}{proved=annex:p03-w03-p03-038-pointed-ef}{The finite game theorem with parameters}
Let \(\tau\) be a finite relational signature, \(k,n\in\N\),
\(\bar a\in(\dom\mathcal A)^n\) and \(\bar b\in(\dom\mathcal B)^n\). The
following are equivalent.
\begin{enumerate}
\item Duplicator has a winning strategy in
      \(G_k(\mathcal A,\bar a;\mathcal B,\bar b)\).
\item \(\operatorname{rec}_k(\mathcal A,\bar a)=
       \operatorname{rec}_k(\mathcal B,\bar b)\).
\item \((\mathcal A,\bar a)\equiv_k(\mathcal B,\bar b)\).
\end{enumerate}
For \(n=0\) this is \cref{thm:p03-ef}. The implications from (2) to (3)
and from (3) to (2) use no choice principle; the passage from (2) to (1)
uses the Axiom of Choice to fix one answer for each Spoiler move when the
structures are infinite.
\end{hbtheorem}

\begin{hbcorollary}{thm:p03-w03-p03-038-finitely-many}{proved=annex:p03-w03-p03-038-finitely-many}{Finitely many formulas of each rank}
For each \(n\) and \(k\), every formula of quantifier rank at most \(k\)
with free variables among \(x_1,\ldots,x_n\) is logically equivalent to
\(\bigvee_{t\in X}\chi_t\) for some \(X\subseteq\mathrm{Rec}_{n,k}\).
Hence there are at most \(2^{|\mathrm{Rec}_{n,k}|}\) such formulas up to
logical equivalence, and in particular finitely many sentences of rank at
most \(k\) up to logical equivalence.
\end{hbcorollary}

\begin{hbcorollary}{thm:p03-w03-p03-038-ee}{proved=annex:p03-w03-p03-038-ee}{Elementary equivalence as all finite games}
For a finite relational signature, \(\mathcal A\equiv\mathcal B\) iff
Duplicator has a winning strategy in \(G_k(\mathcal A,\mathcal B)\) for
every \(k\in\N\).
\end{hbcorollary}

\begin{hbcorollary}{thm:p03-w03-p03-038-criterion}{proved=annex:p03-w03-p03-038-criterion}{The game criterion, proved}
Let \(\mathcal C\) be a class of finite structures in a finite relational
signature. If for every \(k\) there are finite
\(\mathcal A_k\in\mathcal C\) and \(\mathcal B_k\notin\mathcal C\) such that
Duplicator wins \(G_k(\mathcal A_k,\mathcal B_k)\), then no first-order
sentence \(\varphi\) satisfies, for all finite structures \(\mathcal M\),
\(\mathcal M\models\varphi\) iff \(\mathcal M\in\mathcal C\).
\end{hbcorollary}

This is the statement of \cref{thm:p03-ef-inexpressibility}, now with a
proof; its isomorphism-closure hypothesis is not used. The parity example
\cref{def:p03-ef-parity} supplies its own Duplicator strategy for each \(k\)
and so becomes a proved application through this corollary. The criterion
needs one pair for each \(k\): a single pair, or pairs for only some ranks,
would show only that sentences of those ranks fail.

\subsubsection{Connectivity with an explicit round bound}

\begin{hbdefinition}{def:p03-w03-p03-038-cycles}{Cycles, distance and the distance invariant}
For \(N\ge3\), the cycle \(C_N\) has vertex set
\(\mathbb Z/N\mathbb Z\) and edges \(\{u,u+1\}\); the graph
\(C_m\sqcup C_m\) has vertex set
\((\mathbb Z/m\mathbb Z)\times\{0,1\}\) and edges
\(\{(u,e),(u+1,e)\}\). Graphs are structures for one binary symbol \(E\),
interpreted symmetrically and irreflexively. The distance \(d(u,v)\) is
the length of a shortest path, and \(d(u,v)=\infty\) when there is none;
\(\infty\) exceeds every integer. For \(r\in\N\) write
\(x\approx_r y\) when either \(x=y\le2^r\), or \(x>2^r\) and \(y>2^r\).

A position \((\bar a,\bar b)\) of length \(p\) in
\(G_k(C_{2m},C_m\sqcup C_m)\) satisfies the \emph{distance invariant}
\(\mathrm I_r\) when \(d(a_i,a_j)\approx_r d(b_i,b_j)\) for all
\(i,j\le p\). Put \(m_k\defeq(k+1)2^{k+2}\).
\end{hbdefinition}

\begin{hblemma}{thm:p03-w03-p03-038-cycle-response}{proved=annex:p03-w03-p03-038-cycle-response}{Response lemma for cycles}
Let \(k\ge1\), \(m\ge m_k\), \(0\le r<k\), and let a position
\((\bar a,\bar b)\) of length \(p\le k-r-1\) in \(C_{2m}\) and
\(C_m\sqcup C_m\) satisfy \(\mathrm I_{r+1}\). For every vertex \(a\) of
\(C_{2m}\) there is a vertex \(b\) of \(C_m\sqcup C_m\) such that
\((\bar aa,\bar bb)\) satisfies \(\mathrm I_r\); and for every vertex
\(b\) of \(C_m\sqcup C_m\) there is such an \(a\).
\end{hblemma}

\begin{hbtheorem}{thm:p03-w03-p03-038-connectivity}{proved=annex:p03-w03-p03-038-connectivity}{Connectivity is not first-order definable on finite graphs}
For every \(k\in\N\) and every \(m\ge m_k=(k+1)2^{k+2}\), Duplicator has a
winning strategy in \(G_k(C_{2m},C_m\sqcup C_m)\): with \(r\) rounds still
to play, she keeps the distance invariant \(\mathrm I_r\), using
\cref{thm:p03-w03-p03-038-cycle-response}. Consequently \(C_{2m_k}\) and
\(C_{m_k}\sqcup C_{m_k}\) agree through rank \(k\), and no first-order
sentence in the signature \(\{E\}\) is true in exactly the connected finite
graphs.
\end{hbtheorem}

This proves the claims recorded in \cref{def:p03-ef-connectivity}. The bound
\(m_k\) has two roles in the proof: \(m\ge2^{k+2}\) makes every distance
used at most half the length of each cycle, so that it is measured along a
single arc, and \(2m>k(2^k+1)\) leaves a vertex outside all balls of radius
\(2^r\) around previously selected vertices.

\subsubsection{Where finiteness of the signature enters}

\begin{hbwarning}{def:p03-w03-p03-038-infinite-signature}{The converse fails for infinite signatures}
The implication from a winning strategy to agreement on formulas of rank at
most \(k\) holds for any relational signature: the annex proofs of the
implications from (1) to (2) and from (2) to (3) in
\cref{thm:p03-w03-p03-038-pointed-ef} never use the finiteness of
\(\mathrm{At}_n\). The converse needs
\(\mathrm{At}_n\) to be finite. Let \(\tau=\{P_i:i\in\N\}\) with each
\(P_i\) unary, let \(\mathcal A\) have domain \(\N\) with
\(P_i^{\mathcal A}=\{i\}\), and let \(\mathcal B\) extend \(\mathcal A\)
by one element \(e\) in no \(P_i\). Then \(\mathcal A\equiv\mathcal B\),
yet Spoiler wins \(G_1(\mathcal A,\mathcal B)\) by choosing \(e\)
(\cref{thm:p03-w03-p03-038-infinite-signature}).
\end{hbwarning}

\begin{hblemma}{thm:p03-w03-p03-038-infinite-signature}{proved=annex:p03-w03-p03-038-infinite-signature}{The counterexample of the warning}
In \cref{def:p03-w03-p03-038-infinite-signature}, \(\mathcal A\equiv\mathcal B\)
and Spoiler has a winning strategy in \(G_1(\mathcal A,\mathcal B)\).
\end{hblemma}

\subsection{Back-and-forth systems and elementary equivalence}
\label{sec:p03-w03-p03-038-back-and-forth}

A strategy is a dynamic object; a back-and-forth system records the same
information as a static family of finite partial isomorphisms, graded by
the number of rounds that remain.

\begin{hbdefinition}{def:p03-w03-p03-038-bf-system}{Finite back-and-forth systems}
Let \(\mathcal A,\mathcal B\) be structures in a relational signature and
\(k\in\N\). A \emph{\(k\)-back-and-forth system} is a sequence
\((I_j)_{j\le k}\) of sets of positions (of various lengths), each a partial
isomorphism, such that for \(j<k\) and every \((\bar a,\bar b)\in I_{j+1}\):
\begin{enumerate}
\item (forth) for every \(c\in\dom\mathcal A\) there is
      \(d\in\dom\mathcal B\) with \((\bar ac,\bar bd)\in I_j\);
\item (back) for every \(d\in\dom\mathcal B\) there is
      \(c\in\dom\mathcal A\) with \((\bar ac,\bar bd)\in I_j\).
\end{enumerate}
A \emph{back-and-forth family} is a nonempty set \(I\) of partial
isomorphisms with both properties for \(I_{j+1}=I_j=I\). With positions read as
finite maps, this is the potential isomorphism of
\cref{def:p08-bvb501-potential-isomorphism}, since a partial isomorphism
position determines an injective map and conversely.
\end{hbdefinition}

\begin{hbtheorem}{thm:p03-w03-p03-038-bf-system}{proved=annex:p03-w03-p03-038-bf-system}{Systems, games and records}
Let the signature be relational, \(k\in\N\), and let
\((\bar a,\bar b)\) be a position.
\begin{enumerate}
\item Duplicator has a winning strategy in
      \(G_k(\mathcal A,\bar a;\mathcal B,\bar b)\) iff some
      \(k\)-back-and-forth system has \((\bar a,\bar b)\in I_k\).
\item If the signature is finite, the sets
      \[
       I_j\defeq\{(\bar a',\bar b'):\operatorname{rec}_j(\mathcal A,\bar a')
       =\operatorname{rec}_j(\mathcal B,\bar b')\},\qquad j\le k,
      \]
      form a \(k\)-back-and-forth system, and it contains every other one
      coordinatewise.
\item For a finite relational signature, \(\mathcal A\equiv\mathcal B\)
      iff for every \(k\) some \(k\)-back-and-forth system has the empty
      position in \(I_k\).
\end{enumerate}
\end{hbtheorem}

\subsubsection{Bridge from Fra\"iss\'e 1954}
\label{sec:p03-w03-p03-038-fraisse-bridge}

Fra\"iss\'e's 1954 article works with a single \(n\)-ary relation \(R\)
on a base set \(E\), a function from \(E^n\) to \(\{+,-\}\)
(\cref{def:p03-fraisse54-relation}). Its isomorphisms are injections that
preserve the value of every \(n\)-tuple, including tuples with repeated
entries. Read as a structure for one \(n\)-ary relation symbol, an
isomorphism between finite restrictions of \(R\) and of \(S\) is exactly a
finite partial isomorphism in the sense of
\cref{def:p03-w03-p03-038-positions}. The class \(\gamma_R\) of finite
restrictions of \(R\) and their isomorphic copies is the age of \(R\).
Condition \(\mathcal H\) of \cite[\S5.2, p.~377]{Fraisse1954}, restated in
\cref{def:p03-fraisse54-extension}, asks that whenever a finite member of
\(\gamma_R\) extends another, every isomorphism of the smaller onto a
restriction of \(R\) extends to the larger; by \S5.2.1 one added element
suffices.

The article uses this condition in two ways. Theorem~V
\cite[\S5.1, p.~377]{Fraisse1954} states that a homogeneous relation has
every relation of the class \(\Gamma_R\) with finite or countable base
embedded in it. Its proof combines \S5.3 (homogeneity implies
\(\mathcal H\)) with \S5.4(1\textsuperscript{o}), which enumerates the
countable base and extends an isomorphism one element at a time: a forth
construction only. Statement \S5.4(2\textsuperscript{o})
\cite[pp.~378--379]{Fraisse1954} alternates the two bases: if \(R\) and
\(S\) both satisfy \(\mathcal H\), each has every finite restriction of the
other up to isomorphism, and both bases are finite or countable, then
every isomorphism between finite restrictions extends to an isomorphism of
\(R\) onto \(S\). Statement \S5.5 then derives homogeneity from
\(\mathcal H\) on finite or countable bases.

\begin{hbtheorem}{thm:p03-w03-p03-038-fraisse-bridge}{proved=annex:p03-w03-p03-038-fraisse-bridge}{Fra\"iss\'e's extension condition gives the comparison games}
Let \(R\) and \(S\) be \(n\)-ary relations, of any cardinality, that both
satisfy condition \(\mathcal H\) and have the same finite restrictions up
to isomorphism (\(\gamma_R=\gamma_S\)). Then the set of all finite partial
isomorphisms between them is a back-and-forth family. Consequently
Duplicator wins \(G_k(R,\bar a;S,\bar b)\) for every \(k\) and every
partial isomorphism \((\bar a,\bar b)\); in particular \(R\equiv S\). If
both bases are finite or countable, a play of length \(\omega\) in which
Spoiler lists both bases converts any initial partial isomorphism into an
isomorphism of \(R\) onto \(S\), which is \S5.4(2\textsuperscript{o}).
\end{hbtheorem}

\begin{hbexample}{def:p03-w03-p03-038-fraisse-example}{A worked reconstruction: rational and real orders}
Code \((\mathbb Q,<)\) and \((\mathbb R,<)\) as binary relations. Both
have all finite linear orders as their finite restrictions, and both
satisfy \(\mathcal H\): an order embedding of a finite order \(A\) extends
to a one-point extension of \(A\) because every open interval and both end
rays of a dense order without endpoints are nonempty
(\cref{thm:p03-rn-dlo-extension}). By
\cref{thm:p03-w03-p03-038-fraisse-bridge}, Duplicator wins every finite
game between them and the two orders are elementarily equivalent. The
countable-base conclusion \S5.4(2\textsuperscript{o}) does not apply,
since \(\mathbb R\) is uncountable, and it cannot be extended: the two
orders have different cardinalities and are not isomorphic. Thus the
finite back-and-forth content of \(\mathcal H\) survives without the
countability hypothesis, while its isomorphism conclusion does not.
\end{hbexample}

\emph{Source boundary.} The 1954 article contains no game, no quantifier
rank, and no characterization of elementary equivalence. Its only use of
first-order logic is \S10.1 \cite[p.~388]{Fraisse1954}: using a theorem of
Henkin, it derives from Theorems~III and~IV that every arithmetic type in
Tarski's sense, meaning the class of relations satisfying the same
first-order formulas, satisfies the conclusion of Theorem~IV on common
extensions. The proved form of that conclusion in this book is
\cref{thm:p03-fraisse54-logical-joint}. The connection runs only in the
host's direction: by \cref{thm:p03-w03-p03-038-ee}, the arithmetic types of
\S10.1, for one \(n\)-ary relation, are exactly the classes of relations
between any two of which Duplicator wins every finite game. That
characterization, the rank records, and the finite back-and-forth systems
of this subsection are the host's development and are not attributed to
the 1954 article. The article's own route is the forth and alternating
extension of \S5.4 under \(\mathcal H\), together with its filter
constructions; its hypotheses are a single relation of fixed finite arity,
finite or countable bases for \S5.4 and \S5.5, and the Axiom of Choice for
the filter constructions.

\subsubsection{Ehrenfeucht 1961: the games \(G_n\) and \(H_n\)}
\label{sec:p03-w03-p03-038-ehrenfeucht}

Ehrenfeucht's article \cite{Ehrenfeucht1961} introduces comparison games,
games ``in the ordinary sense of the word'' between players I and~II
\cite[\S I.2, p.~130]{Ehrenfeucht1961}, and uses them for two kinds of
indistinguishability: by first-order formulas about individuals, and by
formulas that may also speak about finite sets of individuals. It calls the
first-order criterion a new formulation of a condition of Fra\"iss\'e, citing
Fra\"iss\'e's 1954 publication in the Publications Scientifiques de
l'Universit\'e d'Alger (its reference~[3], p.~141), not the 1954 Annales
article of \cref{sec:p03-w03-p03-038-fraisse-bridge}. The notation below is
the source's: Player~I is Spoiler and Player~II is Duplicator.

\begin{hbdefinition}{def:p03-w03-p03-038-e61-setting}{The languages \(\mathfrak H(A)\) and \(\mathfrak G(A)\)}
Let \(A\) be a set of predicates containing a one-place predicate \(\eta\).
\(\mathfrak H(A)\) is the set of formulas of the lower predicate calculus
with identity \(\iota\) built from a binary predicate \(\epsilon\) and the
predicates of \(A\). A \emph{model} \(\mathfrak M\) for \(\mathfrak H(A)\)
satisfies \cite[\S I.1, p.~129]{Ehrenfeucht1961}:
(1) the elements satisfying \(\eta\), written \(\hat x\eta(x)\), form a set
of individuals; (2) the domain \(|\mathfrak M|\) is the least set \(X\)
with \(\hat x\eta(x)\subseteq X\) that contains \(\{x_1,\ldots,x_k\}\)
whenever \(x_1,\ldots,x_k\in X\), \(k=1,2,\ldots\); (3) the predicates of
\(A\) are interpreted as relations on \(\hat x\eta(x)\); (4) \(\iota\) is
identity; (5) \(\epsilon\) is membership on \(|\mathfrak M|\).
\(\operatorname{stsf}_{\mathfrak M}\alpha\) says that \(\alpha\) is
satisfied in \(\mathfrak M\).

\(\mathfrak G(A)\) is the set of formulas of \(\mathfrak H(A)\) without
\(\epsilon\) whose quantifiers are all restricted to \(\eta\).
\(\mathfrak H_n(A)\) is the set of formulas
\[
(q_kx_k)(q_{k-1}x_{k-1})\cdots(q_1x_1)\Psi(x_1,\ldots,x_n),
\]
each \(q_i\) a quantifier and \(\Psi\) quantifier-free, and
\(\mathfrak G_n(A)\defeq\mathfrak G(A)\cap\mathfrak H_n(A)\).
Two models are \emph{indiscernible by means of finite sets} if they satisfy
the same closed formulas of \(\mathfrak H(A)\) (Definition~1), and
\emph{elementarily indiscernible} if they satisfy the same closed formulas of
\(\mathfrak G(A)\) (Definition~2) \cite[pp.~129--130]{Ehrenfeucht1961}.
\index{Ehrenfeucht--Fraisse game@Ehrenfeucht--Fra\"iss\'e game!Ehrenfeucht's games G_n and H_n}
\end{hbdefinition}

\begin{hbconvention}{def:p03-w03-p03-038-e61-reading}{How the definitions are read}
The source leaves four points implicit; the reconstruction fixes them as
follows, each forced by a use the source makes.
\begin{enumerate}
\item \emph{Individuals are atoms}: no individual is one of the finite sets
      of clause~(2), and no element of \(|\mathfrak M|\) is an
      \(\epsilon\)-member of an individual. The source's support function
      gives an individual \(x\) the support \(\{x\}\), and its claim that a
      matching of individuals determines an isomorphism for \(\epsilon\)
      fails without this.
\item There is at least one individual, as in the proof of Theorem~4.
\item \(\iota\in A\), as in both of the source's examples
      \cite[pp.~131, 134]{Ehrenfeucht1961}; so the winning condition below
      respects identity.
\item In \(\mathfrak G(A)\) a restricted quantifier \((\exists u)_\eta\),
      \((\forall u)_\eta\) is one prefix symbol, abbreviating
      \(\exists u(\eta(u)\wedge\cdot)\) and \(\forall u(\eta(u)\to\cdot)\).
      A formula of \(\mathfrak H_n(A)\) or \(\mathfrak G_n(A)\) \emph{with
      \(k\) free variables} has \(n-k\) prefix quantifiers binding distinct
      variables other than the free ones; this is the count used in the
      proofs of Theorems~1 and~5, one move for each quantifier.
\end{enumerate}
\end{hbconvention}

\begin{hbdefinition}{def:p03-w03-p03-038-e61-games}{The games \(H_n\) and \(G_n\)}
Let \(\mathfrak M_1,\mathfrak M_2\) be models for \(\mathfrak H(A)\). The
game \(H_n(\mathfrak M_1,\mathfrak M_2)\) has \(n\) moves
\cite[\S I.2, p.~130]{Ehrenfeucht1961}. In the \(i\)th move, I chooses
\(l_i\in\{1,2\}\) and a finite sequence
\(a^{l_i}_{i1},\ldots,a^{l_i}_{ik_i}\) of individuals of
\(\mathfrak M_{l_i}\); II answers with \(k_i\) individuals
\(a^{3-l_i}_{i1},\ldots,a^{3-l_i}_{ik_i}\) of \(\mathfrak M_{3-l_i}\). After
\(n\) moves there are \(k_1+\cdots+k_n\) pairs
\(a^1_{ij}\leftrightarrow a^2_{ij}\). II wins iff this correspondence is an
isomorphism for the relations of \(A\): for every \(\alpha\in A\) and every
choice of pairs, \(\alpha\) holds of the chosen entries in \(\mathfrak M_1\)
iff it holds of their partners in \(\mathfrak M_2\). In
\(G_n(\mathfrak M_1,\mathfrak M_2)\), I chooses exactly one individual in
each move; otherwise the rules are those of \(H_n\). A strategy for a player
gives that player's choice in each move as a function of the previous pairs;
it is \emph{winning} if the player wins against every strategy of the other
\cite[pp.~131--132]{Ehrenfeucht1961}. ``After \(k\) moves'' a strategy is
taken from a given position of \(k\) moves.
\end{hbdefinition}

\begin{hblemma}{thm:p03-w03-p03-038-e61-determined}{proved=annex:p03-w03-p03-038-e61-determined}{The games are determined}
From every position of \(H_n(\mathfrak M_1,\mathfrak M_2)\) or
\(G_n(\mathfrak M_1,\mathfrak M_2)\) exactly one player has a winning
strategy. The proof uses the Axiom of Choice when the models are infinite.
\end{hblemma}

The source records determinacy as obvious \cite[p.~132]{Ehrenfeucht1961}; the
annex proves it by induction on the number of remaining moves.

\begin{hbexample}{def:p03-w03-p03-038-e61-int-rat}{Integers and rationals}
Let \(A\) consist of \(\eta\), \(\iota\) and a binary \(\varrho\), and let
the individuals of \(\mathfrak M_1\) and \(\mathfrak M_2\) be the integers
and the rationals, with \(\varrho\) read as \(<\)
\cite[p.~131]{Ehrenfeucht1961}. II has a winning strategy in
\(H_1(\mathfrak M_1,\mathfrak M_2)\), while I has a winning strategy in
\(G_3(\mathfrak M_1,\mathfrak M_2)\): he chooses \(0\) and \(1\) among the
integers and then the midpoint of II's two answers
(\cref{thm:p03-w03-p03-038-e61-int-rat}).
\end{hbexample}

\begin{hblemma}{thm:p03-w03-p03-038-e61-int-rat}{proved=annex:p03-w03-p03-038-e61-int-rat}{The strategies of the example}
In \cref{def:p03-w03-p03-038-e61-int-rat}, II wins
\(H_1(\mathfrak M_1,\mathfrak M_2)\) and I wins
\(G_3(\mathfrak M_1,\mathfrak M_2)\).
\end{hblemma}

\begin{hbdefinition}{def:p03-w03-p03-038-e61-support}{Supports and established isomorphisms}
For \(x\in|\mathfrak M|\) let \(S(x)\defeq\{x\}\) if \(x\) is an
individual, and otherwise let \(S(x)\) be the least set \(Y\) with
\(x\in Y\) that contains every member of each of its elements; put
\(s(x)\defeq S(x)\cap\hat x\eta(x)\) \cite[p.~132]{Ehrenfeucht1961}. For a
set \(X\) of individuals, \(\operatorname{HF}(X)\) is the least set
containing \(X\) and closed under forming
\(\{x_1,\ldots,x_k\}\), \(k\ge1\). A correspondence
\(a_1\leftrightarrow b_1,\ldots,a_m\leftrightarrow b_m\) of individuals that
is an isomorphism for \(A\) \emph{establishes an isomorphism between}
\(x\in|\mathfrak M_1|\) and \(y\in|\mathfrak M_2|\) when
\(s(x)\subseteq\{a_1,\ldots,a_m\}\), \(s(y)\subseteq\{b_1,\ldots,b_m\}\),
and the map \(F\) of \cref{thm:p03-w03-p03-038-e61-hull} sends \(x\) to
\(y\).
\end{hbdefinition}

\begin{hblemma}{thm:p03-w03-p03-038-e61-hull}{proved=annex:p03-w03-p03-038-e61-hull}{A matching of individuals extends to the finite sets over them}
Under \cref{def:p03-w03-p03-038-e61-reading}, let
\(a_1\leftrightarrow b_1,\ldots,a_m\leftrightarrow b_m\) be an isomorphism
for \(A\) between individuals of \(\mathfrak M_1\) and \(\mathfrak M_2\),
and \(X=\{a_1,\ldots,a_m\}\), \(Y=\{b_1,\ldots,b_m\}\). Then
\(\operatorname{HF}(X)=\{x\in|\mathfrak M_1|:s(x)\subseteq X\}\), and there
is exactly one map \(F:\operatorname{HF}(X)\to\operatorname{HF}(Y)\) with
\(F(a_i)=b_i\) and \(F(x)=\{F(t):t\in x\}\) for non-individuals \(x\). It is
a bijection, and for every atomic formula \(\alpha\) of \(\mathfrak H(A)\)
and elements \(z_1,\ldots\) of \(\operatorname{HF}(X)\),
\(\operatorname{stsf}_{\mathfrak M_1}\alpha(z_1,\ldots)\) iff
\(\operatorname{stsf}_{\mathfrak M_2}\alpha(F(z_1),\ldots)\).
\end{hblemma}

\begin{hbtheorem}{thm:p03-w03-p03-038-e61-thm1}{proved=annex:p03-w03-p03-038-e61-thm1}{Theorems 1--3 of the source}
Under \cref{def:p03-w03-p03-038-e61-reading}:
\begin{enumerate}
\item (Theorem~1, p.~132) Let \(x_1,\ldots,x_k\in|\mathfrak M_1|\) and
      \(y_1,\ldots,y_k\in|\mathfrak M_2|\). If II has a winning strategy in
      \(H_n(\mathfrak M_1,\mathfrak M_2)\) after \(k\) moves whose
      correspondence establishes an isomorphism between \(x_i\) and \(y_i\)
      for \(i=1,\ldots,k\), then for every \(\Psi\in\mathfrak H_n(A)\) with
      \(k\) free variables,
      \(\operatorname{stsf}_{\mathfrak M_1}\Psi(x_1,\ldots,x_k)\equiv
      \operatorname{stsf}_{\mathfrak M_2}\Psi(y_1,\ldots,y_k)\).
\item (Theorem~2, p.~133) If II has a winning strategy in
      \(H_n(\mathfrak M_1,\mathfrak M_2)\), the models agree on every closed
      formula of \(\mathfrak H_n(A)\).
\item (Theorem~3, p.~133) If II has a winning strategy in
      \(H_n(\mathfrak M_1,\mathfrak M_2)\) for every \(n\), the models are
      indiscernible by means of finite sets.
\end{enumerate}
\end{hbtheorem}

\begin{hbtheorem}{thm:p03-w03-p03-038-e61-thm4}{proved=annex:p03-w03-p03-038-e61-thm4}{Theorem 4 of the source: a countable arithmetical submodel}
Under \cref{def:p03-w03-p03-038-e61-reading}, let \(A\) be finite or
countable and \(\mathfrak M\) a model for \(\mathfrak H(A)\). Then some
model \(\mathfrak M'\) for \(\mathfrak H(A)\) with countably infinite domain
is an elementary submodel of \(\mathfrak M\) for \(\mathfrak H(A)\). The
proof uses the Axiom of Choice.
\end{hbtheorem}

The source states Theorem~4 without the hypothesis on \(A\); its proof
introduces one Skolem function for each formula and counts \(\aleph_0\) of
them \cite[p.~134]{Ehrenfeucht1961}, which is where the countability of
\(A\) enters. ``Arithmetical submodel'' is the source's term for elementary
submodel.

\begin{hbtheorem}{thm:p03-w03-p03-038-e61-sec3}{proved=annex:p03-w03-p03-038-e61-sec3}{Indiscernible models on which I wins \(H_3\)}
Let \(A\) consist of \(\iota,\eta\) and predicates \(\alpha(x)\),
\(\beta(x)\), \(\gamma(x,y)\), \(\delta(x,y,z)\). Let the individuals of
\(\mathfrak M_1\) be the infinite sequences with values \(0,1\) and the
natural numbers (taken as atoms), with \(\alpha\): ``\(x\) is a \(0\)-\(1\)
sequence'', \(\beta\): ``\(x\) is a natural number'', \(\gamma\): ``\(x,y\)
are natural numbers and \(x<y\)'', \(\delta\): ``\(z\) is \(0\) or \(1\),
\(y\) is a natural number, \(x\) is a sequence and its \(y\)th term is
\(z\)''. Let \(\mathfrak M_2\) be a countable arithmetical submodel of
\(\mathfrak M_1\) (\cref{thm:p03-w03-p03-038-e61-thm4}). Then
\(\mathfrak M_1\) and \(\mathfrak M_2\) are indiscernible by means of finite
sets, yet I has a winning strategy in \(H_3(\mathfrak M_1,\mathfrak M_2)\)
\cite[\S I.3, pp.~134--135]{Ehrenfeucht1961}. So the converse of Theorem~3
fails.
\end{hbtheorem}

\begin{hbtheorem}{thm:p03-w03-p03-038-e61-thm5}{proved=annex:p03-w03-p03-038-e61-thm5}{Theorems 5--7 of the source}
Under \cref{def:p03-w03-p03-038-e61-reading}:
\begin{enumerate}
\item (Theorem~5, p.~135) If II has a winning strategy in
      \(G_n(\mathfrak M_1,\mathfrak M_2)\) after \(k\) moves
      \(x_1\leftrightarrow y_1,\ldots,x_k\leftrightarrow y_k\), then
      \(\operatorname{stsf}_{\mathfrak M_1}\Psi(x_1,\ldots,x_k)\equiv
      \operatorname{stsf}_{\mathfrak M_2}\Psi(y_1,\ldots,y_k)\) for every
      \(\Psi\in\mathfrak G_n(A)\) with \(k\) free variables.
\item (Theorem~6, p.~135) If II has a winning strategy in
      \(G_n(\mathfrak M_1,\mathfrak M_2)\), the models agree on every closed
      formula of \(\mathfrak G_n(A)\).
\item (Theorem~7, p.~136) If II has a winning strategy in
      \(G_n(\mathfrak M_1,\mathfrak M_2)\) for every \(n\), the models are
      elementarily indiscernible.
\end{enumerate}
\end{hbtheorem}

\begin{hbtheorem}{thm:p03-w03-p03-038-e61-converse-fails}{proved=annex:p03-w03-p03-038-e61-converse-fails}{The converse of Theorem 7 fails for infinite \(A\)}
Let \(A\) contain \(\iota\), \(\eta\) and one-place predicates
\(a_1,a_2,\ldots\). Let the individuals of \(\mathfrak M_1\) be the natural
numbers with \(a_i(m)\) iff \(m>i\), and those of \(\mathfrak M_2\) the
ordinals \(\le\omega\) with the same interpretation
\cite[p.~136]{Ehrenfeucht1961}. Then I wins
\(G_1(\mathfrak M_1,\mathfrak M_2)\) by choosing \(\omega\), and the two
models are elementarily indiscernible.
\end{hbtheorem}

\begin{hblemma}{thm:p03-w03-p03-038-unary-types}{proved=annex:p03-w03-p03-038-unary-types}{Counting unary types}
Let \(P\) be a finite set of one-place predicates and \(\mathcal A,\mathcal
B\) structures for \(P\) (with equality). For \(t\subseteq P\) let
\(n_t(\mathcal A)\) be the number, finite or infinite, of elements of
\(\mathcal A\) satisfying exactly the predicates in \(t\). If for every
\(t\) the numbers \(n_t(\mathcal A)\) and \(n_t(\mathcal B)\) are equal or
both infinite, then the finite partial isomorphisms between \(\mathcal A\)
and \(\mathcal B\) form a back-and-forth family containing the empty
position, and \(\mathcal A\equiv\mathcal B\).
\end{hblemma}

\begin{hbdefinition}{def:p03-w03-p03-038-discrete-orders}{Discrete orders with a least element}
A linear order \((L,<)\) satisfies \((\mathrm D)\) if it has a least element
and no greatest element, every element has an immediate successor, and every
element other than the least has an immediate predecessor. Both \(\omega\)
and \(\omega+\omega^*+\omega\) satisfy \((\mathrm D)\).
\end{hbdefinition}

\begin{hblemma}{thm:p03-w03-p03-038-discrete-orders}{proved=annex:p03-w03-p03-038-discrete-orders}{Gap strategy for discrete orders}
If two linear orders satisfy \((\mathrm D)\), Duplicator wins
\(G_k\) between them for every \(k\); hence they are elementarily
equivalent in the signature \(\{<\}\).
\end{hblemma}

\begin{hbtheorem}{thm:p03-w03-p03-038-e61-remark}{proved=annex:p03-w03-p03-038-e61-remark}{The source's remark after Theorem 7}
Under \cref{def:p03-w03-p03-038-e61-reading} \cite[p.~136]{Ehrenfeucht1961}:
\begin{enumerate}
\item a winning strategy for II in \(G_n\) does not yield one in \(H_n\): for
      the models of \cref{thm:p03-w03-p03-038-e61-sec3}, II wins every
      \(G_n\) and I wins \(H_3\);
\item two models can be elementarily indiscernible but discernible by means
      of finite sets: take \(A=\{\iota,\eta,\varrho\}\) with \(\varrho\) an
      ordering of type \(\omega\) in one model and of type
      \(\omega+\omega^*+\omega\) in the other.
\end{enumerate}
\end{hbtheorem}

For the first claim the source points to an application of Theorem~8; the
annex derives it from \cref{thm:p03-w03-p03-038-e61-sec3} and Theorem~10.

\begin{hbdefinition}{def:p03-w03-p03-038-e61-omega}{The formulas \(\Omega_{kl}\)}
Let \(A\) be finite and \(n\in\N\) \cite[pp.~136--137]{Ehrenfeucht1961}.
(i) Let \(\Omega_{n0},\ldots,\Omega_{nl_n}\) be quantifier-free formulas of
\(\mathfrak G(A)\) in the variables \(u_1,\ldots,u_n\) only, such that
\(\Omega_{n0}\vee\cdots\vee\Omega_{nl_n}\) and
\(\sim(\Omega_{ni}\wedge\Omega_{nj})\), \(i\ne j\), are tautologies, and for
every quantifier-free \(\Gamma\) of \(\mathfrak G(A)\) in
\(u_1,\ldots,u_n\) and every \(i\), either \(\Omega_{ni}\supset\Gamma\) or
\(\Omega_{ni}\supset\sim\Gamma\) is a tautology. Such a list exists because
\(A\) is finite: take every conjunction that contains, for each atomic
formula in \(u_1,\ldots,u_n\), either it or its negation. (ii) For
\(1\le k\le n\), let \(\Omega_{k-1,0},\ldots,\Omega_{k-1,l_{k-1}}\) list all
conjunctions \(A_0\wedge\cdots\wedge A_{l_k}\) in which each \(A_i\) is
\((\exists u_k)\Omega_{ki}\) or \(\sim(\exists u_k)\Omega_{ki}\), the
quantifier restricted to \(\eta\).
\end{hbdefinition}

\begin{hbtheorem}{thm:p03-w03-p03-038-e61-thm8}{proved=annex:p03-w03-p03-038-e61-thm8}{Theorems 8--10 of the source}
Let \(A\) be finite and assume \cref{def:p03-w03-p03-038-e61-reading}.
\begin{enumerate}
\item (Lemma, p.~137) If
      \(\operatorname{stsf}_{\mathfrak M_1}\Omega_{ki}(x_1,\ldots,x_k)\), then
      for every \(\mathfrak M_2\) and \(y_1,\ldots,y_k\), II has a winning
      strategy in \(G_n(\mathfrak M_1,\mathfrak M_2)\) after the moves
      \(x_j\leftrightarrow y_j\) iff
      \(\operatorname{stsf}_{\mathfrak M_2}\Omega_{ki}(y_1,\ldots,y_k)\).
\item (Theorem~8, p.~136) For every \(\mathfrak M_1\) and individuals
      \(x_1,\ldots,x_k\) there is \(\Phi\in\mathfrak G(A)\) such that for all
      \(\mathfrak M_2\) and \(y_1,\ldots,y_k\), II has a winning strategy in
      \(G_n(\mathfrak M_1,\mathfrak M_2)\) after the moves
      \(x_j\leftrightarrow y_j\) iff
      \(\operatorname{stsf}_{\mathfrak M_2}\Phi(y_1,\ldots,y_k)\).
\item (Theorem~9, p.~137) For every \(n\) there are closed formulas
      \(\Phi_1,\ldots,\Phi_s\) of \(\mathfrak G(A)\) such that II has a
      winning strategy in \(G_n(\mathfrak M_1,\mathfrak M_2)\) iff
      \(\operatorname{stsf}_{\mathfrak M_1}\Phi_i\equiv
      \operatorname{stsf}_{\mathfrak M_2}\Phi_i\) for every \(i\).
\item (Theorem~10, p.~137) II has a winning strategy in
      \(G_n(\mathfrak M_1,\mathfrak M_2)\) for every \(n\) iff the models are
      elementarily indiscernible.
\end{enumerate}
\end{hbtheorem}

As printed, the lemma says that \(\Omega_{ki}\) satisfies ``the conditions
of Theorem~7'', and Theorem~9 is said to follow from ``Theorem~3 and the
lemma''; the content used is the condition of Theorem~8 and the lemma alone.

\emph{Bridge to this section.} The individuals of a model with the relations
of \(A\setminus\{\iota,\eta\}\) form a relational structure, and
\(G_n(\mathfrak M_1,\mathfrak M_2)\) is the game
\(G_n\) of \cref{def:p03-w03-p03-038-positions} between these structures.
A formula of \(\mathfrak G_n(A)\) with \(k\) free variables has quantifier
rank \(n-k\), so Theorems~5 and~6 are the prenex case of the implication
from (1) to (3) in \cref{thm:p03-w03-p03-038-pointed-ef}. The formulas
\(\Omega_{kl}\) play the part of the characteristic formulas
\(\chi_t\) of \cref{def:p03-w03-p03-038-characteristic}, by a different
construction: \(\chi_t\) is indexed by records and carries a universal
clause, while \(\Omega_{k-1,l}\) runs over all sign patterns of existential
formulas, and the universal clause is implied because the
\(\Omega_{kl}\) are exhaustive. Neither family is prenex. Theorem~10 is
\cref{thm:p03-w03-p03-038-ee}, and
\cref{thm:p03-w03-p03-038-e61-converse-fails} is a second witness, with
different structures, for \cref{def:p03-w03-p03-038-infinite-signature}.
\emph{Source boundary.} The source measures formulas by the number of
quantifiers in a prenex prefix (Theorems~1 and~5) and by the nesting of
its \(\Omega\) formulas; it defines no quantifier rank for arbitrary
formulas, and its games are played on individuals only, also in \(H_n\),
where the formulas range over finite sets. The source has two parts:
Part~I (pp.~129--137), presented above, and Part~II (pp.~138--141), which
applies the games \(H_n\) to formalized theories of ordinal arithmetic and
is presented in \cref{sec:p03-w03-p03-038-e61-ordinals}. Apart from that
application, the games \(H_n\), the support function and Theorems~1--4 have
no counterpart elsewhere in this section.

\subsubsection{Ehrenfeucht 1961, Part~II: ordinal arithmetic}
\label{sec:p03-w03-p03-038-e61-ordinals}

Part~II of \cite{Ehrenfeucht1961} (pp.~138--141) uses the games \(H_n\) to
compare structures of ordinals under order, addition and multiplication.
Its conclusions are that ordinals congruent modulo \(\omega^n\) give
structures on which II wins \(H_n\) (Theorem~12), that three lifting steps
produce structures indiscernible by means of finite sets (Theorems~13--17),
and that an ordinal is definable in the order, additive or multiplicative
theory exactly when it lies below \(\omega^{\omega}\),
\(\omega^{\omega^{\omega}}\) or \(\omega^{\omega^{\omega^{\omega}}}\)
respectively (Theorem~18), which the source presents as a conjecture of
Tarski. \emph{Inputs:} the setting and games of
\cref{def:p03-w03-p03-038-e61-setting,def:p03-w03-p03-038-e61-games}
under \cref{def:p03-w03-p03-038-e61-reading}; Theorems~1--3
(\cref{thm:p03-w03-p03-038-e61-thm1}); ordinal arithmetic as defined in
\cref{def:p06-ordinal-arithmetic}, with the laws of
\cref{thm:p06-ordinal-arithmetic-laws} and the Cantor normal form of
\cref{thm:p06-cantor-normal-form}. The Axiom of Choice enters only where
a winning strategy is extracted from the relation \(W_j\) below.

\begin{hbdefinition}{def:p03-w03-p03-038-e61-ordinal-models}{The theories \(\mathcal T_i\), \(T_i\) and the models \(\mathfrak M_a\)}
Following \cite[p.~138]{Ehrenfeucht1961}, let \(A_1\) consist of a binary
predicate \(\nu\), let \(A_2\) consist of \(\nu\) and a ternary \(\sigma\),
and let \(A_3\) consist of \(\nu\), \(\sigma\) and a ternary \(\mu\); each
\(A_i\) also contains \(\eta\) (by \cref{def:p03-w03-p03-038-e61-setting})
and \(\iota\) (by \cref{def:p03-w03-p03-038-e61-reading}(3)). Put
\(\mathcal T_i\defeq\mathfrak H(A_i)\) and \(T_i\defeq\mathfrak G(A_i)\).
For an ordinal \(a\), \(\mathfrak M_a\) is the model for \(\mathcal T_i\)
whose individuals are the ordinals \(x<a\), with \(\nu(x,y)\) read as
\(x<y\), \(\sigma(x,y,z)\) as \(x+y=z\) and \(\mu(x,y,z)\) as
\(x\cdot y=z\), all arguments being individuals of the model. For
\(a<b\), \(\mathfrak M_{ab}\) is the model whose individuals are the
ordinals \(x\) with \(a<x<b\). The source also uses the \emph{class of
all ordinals} as a model \cite[p.~140]{Ehrenfeucht1961}; we write
\(\mathfrak M_\infty\) for it and \(\mathfrak M_{a\infty}\) for its
individuals \(x>a\). As required by
\cref{def:p03-w03-p03-038-e61-reading}(1), the individuals are atoms,
one for each ordinal, written as the ordinal itself; the domain of
\(\mathfrak M_\infty\) is the class built over them by clause~(2), and its
satisfaction relation is given, one formula at a time, by the recursive
clauses. A model \(\mathfrak M_a\) with \(a=0\), or a segment with no
individuals, is not a model in the sense of clause~(2); it still occurs as
a component of the games below, which involve individuals only.
\index{Ehrenfeucht--Fraisse game@Ehrenfeucht--Fra\"iss\'e game!ordinal arithmetic}
\end{hbdefinition}

\begin{hbdefinition}{def:p03-w03-p03-038-e61-wins}{Winning positions \(W_j\)}
Let \(\mathfrak M_1,\mathfrak M_2\) be structures whose individuals,
possibly none and possibly a proper class, carry the relations of \(A\). A
\emph{position} is a finite set of pairs \(x\leftrightarrow y\) of an
individual \(x\) of \(\mathfrak M_1\) and an individual \(y\) of
\(\mathfrak M_2\). Define \(W_0(P)\) to hold iff \(P\) is an isomorphism
for \(A\) in the sense of \cref{def:p03-w03-p03-038-e61-games}, and
\(W_{j+1}(P)\) to hold iff for every move of I in the sense of \(H_n\)
(a choice of \(l\in\{1,2\}\) and a nonempty finite sequence of individuals
of \(\mathfrak M_l\)) there is an answer of II (a sequence of the same
length in \(\mathfrak M_{3-l}\)) such that \(W_j\) holds of \(P\) together
with the new pairs. Write \(W_j(\mathfrak M_1,\mathfrak M_2)\) for
\(W_j(\emptyset)\). For a fixed \(j\) this is a statement about the two
structures, also when one of them is \(\mathfrak M_\infty\).
\end{hbdefinition}

\begin{hblemma}{thm:p03-w03-p03-038-e61-wins}{proved=annex:p03-w03-p03-038-e61-wins}{Properties of \(W_j\)}
For structures as in \cref{def:p03-w03-p03-038-e61-wins}:
\begin{enumerate}
\item \(W_j(P\cup Q)\) implies \(W_j(P)\); \(W_{j+1}(P)\) implies
      \(W_j(P)\); and \(W_j(P)\) implies that \(P\) is an isomorphism for
      \(A\).
\item If \(\pi_1,\pi_2\) are bijections of the individuals of
      \(\mathfrak M_1,\mathfrak M_2\) onto those of
      \(\mathfrak M_1',\mathfrak M_2'\) that are isomorphisms for \(A\),
      then \(W_j(P)\) iff \(W_j\) holds of the image of \(P\).
\item If the individuals of both structures form sets, II has a winning
      strategy in \(H_n(\mathfrak M_1,\mathfrak M_2)\) after a position
      \(P\) of \(k\) moves iff \(W_{n-k}(P)\). The direction from
      \(W_{n-k}(P)\) to a strategy uses the Axiom of Choice.
\item (Theorems~1--3 in terms of \(W_j\).) Let \(\mathfrak M_1,
      \mathfrak M_2\) be models for \(\mathfrak H(A)\) under
      \cref{def:p03-w03-p03-038-e61-reading}, whose individuals may form
      a proper class. If \(W_j(P)\), \(P\) establishes an isomorphism
      between \(x_i\) and \(y_i\) for \(i\le k\), and
      \(\Psi\in\mathfrak H(A)\) is prenex with at most \(j\) prefix
      quantifiers and free variables among the first \(k\), then
      \(\operatorname{stsf}_{\mathfrak M_1}\Psi(x_1,\ldots,x_k)\equiv
      \operatorname{stsf}_{\mathfrak M_2}\Psi(y_1,\ldots,y_k)\). If
      \(W_n(\mathfrak M_1,\mathfrak M_2)\) holds for every \(n\), the two
      models are indiscernible by means of finite sets, hence elementarily
      indiscernible.
\end{enumerate}
\end{hblemma}

\begin{hbdefinition}{def:p03-w03-p03-038-e61-congruence}{Congruence of ordinals}
For ordinals \(a,b\) and \(c>0\), \(a\equiv b\pmod c\) iff either \(a=b\)
or \(a=c\cdot d+e\) and \(b=c\cdot d'+e\) for some ordinals \(e\) and
\(d,d'\ne0\) \cite[p.~138]{Ehrenfeucht1961}. The source adds a side
condition on \(c\), phrased in terms of cofinality; the moduli to which it
applies the definition are \(\omega^n\) and \(\omega^\omega\), and the
arguments below use only the two displayed equations. For \(x<y\), the
\emph{type} of the open segment \(\{z:x<z<y\}\) is the unique \(t\) with
\(x+1+t=y\) (\cref{thm:p03-w03-p03-038-e61-toolkit}(2)); the map
\(t'\mapsto x+1+t'\) is an order isomorphism from \(t\) onto that
segment, and \(t'\mapsto x+1+t'\) maps the class of all ordinals onto the
individuals of \(\mathfrak M_{x\infty}\). The function
\(\operatorname{sgn}x\) is \(0\) for \(x=0\) and \(1\) for \(x>0\).
\end{hbdefinition}

\begin{hblemma}{thm:p03-w03-p03-038-e61-toolkit}{proved=annex:p03-w03-p03-038-e61-toolkit}{Ordinal arithmetic used in Part~II}
For ordinals \(\alpha,\beta,\xi,v,x,y\):
\begin{enumerate}
\item \(\xi\le\alpha+\xi\); if \(\alpha\ge1\), \(\xi\le\alpha\cdot\xi\);
      if \(\alpha\le\alpha'\), then \(\alpha+\beta\le\alpha'+\beta\) and
      \(\alpha\cdot\beta\le\alpha'\cdot\beta\).
\item If \(\alpha\le\beta\), there is exactly one \(\gamma\) with
      \(\alpha+\gamma=\beta\).
\item If \(\delta>0\), there are unique \(q\) and \(r<\delta\) with
      \(v=\delta\cdot q+r\).
\item If \(x=\omega^{e_1}m_1+\cdots+\omega^{e_r}m_r\) is a Cantor normal
      form with \(r\ge1\), then
      \(\omega^{e_1}m_1\le x<\omega^{e_1}(m_1+1)\le\omega^{e_1+1}\).
      Hence \(x<\omega^\gamma\) iff \(x=0\) or \(e_1<\gamma\), and the
      ordinals below \(\omega^\gamma\) are exactly the Cantor normal forms
      whose exponents are below \(\gamma\).
\item If \(x,y<\omega^\gamma\), then \(x+y<\omega^\gamma\); if
      \(\gamma>0\), \(\omega^\gamma\) is a limit ordinal; if
      \(x<\omega^\gamma\) and \(c\ge1\), then
      \(x+\omega^\gamma c=\omega^\gamma c\).
\item (Comparison.) For Cantor normal forms \(x\) and \(y\), \(x<y\) iff
      the list of terms of \(x\) is a proper initial part of that of
      \(y\), or at the first place where the lists differ the exponent of
      \(x\) is smaller, or the exponents agree and the coefficient of
      \(x\) is smaller.
\item (Addition.) If \(y\ne0\) has leading term \(\omega^hp_1\), then
      \(x+y\) is obtained from \(x\) by deleting its terms of exponent
      below \(h\), adding \(p_1\) to the coefficient of \(\omega^h\) (or
      inserting \(\omega^hp_1\) when \(x\) has no such term), and
      appending the remaining terms of \(y\).
\item (Multiplication.) If \(x,y\ne0\), \(x\) has leading exponent
      \(e_1\), leading coefficient \(m_1\) and remaining terms \(x'\), and
      \(y=\sum_{i\le s}\omega^{h_i}p_i\), then
      \(x\cdot y=\sum_{i:h_i>0}\omega^{e_1+h_i}p_i\), followed, when
      \(h_s=0\), by \(\omega^{e_1}(m_1p_s)+x'\); this is a Cantor normal
      form.
\end{enumerate}
\end{hblemma}

\begin{hblemma}{thm:p03-w03-p03-038-e61-sums}{proved=annex:p03-w03-p03-038-e61-sums}{Theorem~11 and the segment lemma}
Let the structures carry \(\eta\), \(\iota\) and an order \(\nu\), as for
\(\mathcal T_1\).
\begin{enumerate}
\item (Ordered sums.) Let \(L=L_0\oplus\cdots\oplus L_r\) and
      \(L'=L_0'\oplus\cdots\oplus L_r'\) be ordered sums of linear orders,
      each a set or a class, possibly empty. If a position \(P\) pairs
      elements of \(L_i\) only with elements of \(L_i'\), and its part
      \(P_i\) in \(L_i\times L_i'\) satisfies \(W_j(P_i)\) for every
      \(i\), then \(W_j(P)\) holds for \((L,L')\).
\item (Theorem~11, p.~138.) Let \(a_1<\cdots<a_k<a\) and
      \(b_1<\cdots<b_k<b\), where \(b\) may be \(\infty\). If
      \(W_{n-1}\) holds for \((\mathfrak M_{a_1},\mathfrak M_{b_1})\), for
      \((\mathfrak M_{a_ia_{i+1}},\mathfrak M_{b_ib_{i+1}})\),
      \(1\le i<k\), and for \((\mathfrak M_{a_ka},\mathfrak M_{b_kb})\),
      then II wins \(H_n(\mathfrak M_a,\mathfrak M_b)\) after the first
      move \(a_1\leftrightarrow b_1,\ldots,a_k\leftrightarrow b_k\), in the
      sense that \(W_{n-1}\) holds of this position.
\item (Lemma, p.~139.) If \(W_n(\mathfrak M_a,\mathfrak M_b)\), then
      \(W_n(\mathfrak M_{a+d},\mathfrak M_{b+d})\) for every ordinal
      \(d\).
\end{enumerate}
\end{hblemma}

\begin{hbtheorem}{thm:p03-w03-p03-038-e61-thm12}{proved=annex:p03-w03-p03-038-e61-thm12}{Theorem~12 of the source and its class form}
Let \(n\ge1\), with models for \(\mathcal T_1\).
\begin{enumerate}
\item (Theorem~12, p.~138.) If \(a\equiv b\pmod{\omega^n}\), then
      \(W_n(\mathfrak M_a,\mathfrak M_b)\); so II has a winning strategy in
      \(H_n(\mathfrak M_a,\mathfrak M_b)\).
\item If \(a=\omega^n\cdot c\) with \(c\ge1\), then
      \(W_n(\mathfrak M_a,\mathfrak M_\infty)\).
\end{enumerate}
\end{hbtheorem}

\begin{hbtheorem}{thm:p03-w03-p03-038-e61-thm13}{proved=annex:p03-w03-p03-038-e61-thm13}{Theorem~13 of the source}
If \(a\equiv b\pmod{\omega^\omega}\), then \(\mathfrak M_a\) and
\(\mathfrak M_b\), as models for \(\mathcal T_1\), are indiscernible by
means of finite sets \cite[p.~139]{Ehrenfeucht1961}, and so elementarily
indiscernible. For every \(c\ge1\), \(\mathfrak M_{\omega^\omega c}\) and
\(\mathfrak M_\infty\) are indiscernible by means of finite sets as models
for \(\mathcal T_1\).
\end{hbtheorem}

For Theorems~14 and~16, let \(f\) be a finite injective map from a set of
ordinals to ordinals that preserves \(<\). Let \(E(f)\) be the set of
ordinals whose Cantor normal form has all exponents in the domain of
\(f\), and for \(x=\sum_j\omega^{e_j}m_j\in E(f)\) put
\(g_f(x)\defeq\sum_j\omega^{f(e_j)}m_j\), with \(g_f(0)=0\).

\begin{hblemma}{thm:p03-w03-p03-038-e61-lift}{proved=annex:p03-w03-p03-038-e61-lift}{Lifting a correspondence of exponents}
\begin{enumerate}
\item \(g_f\) is a bijection of \(E(f)\) onto \(E(f^{-1})\) with inverse
      \(g_{f^{-1}}\), and for \(x,y,z\in E(f)\): \(x<y\) iff
      \(g_f(x)<g_f(y)\), and \(x+y=z\) iff \(g_f(x)+g_f(y)=g_f(z)\).
\item If moreover \(u+v=w\) iff \(f(u)+f(v)=f(w)\) for \(u,v,w\) in the
      domain of \(f\), then \(x\cdot y=z\) iff
      \(g_f(x)\cdot g_f(y)=g_f(z)\) for \(x,y,z\in E(f)\).
\item The least ordinal greater than every
      \(\omega^{a_1}n_1+\cdots+\omega^{a_k}n_k\) with \(a_1,\ldots,a_k<a\)
      and \(n_1,\ldots,n_k<\omega\) is \(\omega^a\).
\end{enumerate}
\end{hblemma}

\begin{hbtheorem}{thm:p03-w03-p03-038-e61-thm14}{proved=annex:p03-w03-p03-038-e61-thm14}{Theorems~14 and~16 of the source}
Let \(c\) be the least ordinal greater than every
\(\omega^{a_1}n_1+\cdots+\omega^{a_k}n_k\) with \(a_i<a\),
\(n_i<\omega\), and \(d\) the corresponding ordinal for \(b\); by
\cref{thm:p03-w03-p03-038-e61-lift}(3), \(c=\omega^a\) and
\(d=\omega^b\).
\begin{enumerate}
\item (Theorem~14, p.~139.) If \(W_n(\mathfrak M_a,\mathfrak M_b)\) for
      \(\mathcal T_1\), then \(W_n(\mathfrak M_c,\mathfrak M_d)\) for
      \(\mathcal T_2\): II has a winning method in
      \(H_n(\mathfrak M_c,\mathfrak M_d)\).
\item (Theorem~16, p.~140.) If \(W_n(\mathfrak M_a,\mathfrak M_b)\) for
      \(\mathcal T_2\), then \(W_n(\mathfrak M_c,\mathfrak M_d)\) for
      \(\mathcal T_3\).
\item Both statements hold with \(b\) and \(d\) replaced by \(\infty\).
\end{enumerate}
\end{hbtheorem}

\begin{hbtheorem}{thm:p03-w03-p03-038-e61-thm15}{proved=annex:p03-w03-p03-038-e61-thm15}{Theorems~15 and~17 of the source}
Let \(a\equiv b\pmod{\omega^\omega}\).
\begin{enumerate}
\item (Theorem~15, p.~140.) \(\mathfrak M_{\omega^a}\) and
      \(\mathfrak M_{\omega^b}\), as models for \(\mathcal T_2\), are
      indiscernible by means of finite sets.
\item (Theorem~17, p.~140.) \(\mathfrak M_{\omega^{\omega^a}}\) and
      \(\mathfrak M_{\omega^{\omega^b}}\), as models for \(\mathcal T_3\),
      are indiscernible by means of finite sets.
\item \(\mathfrak M_{\omega^{\omega^\omega}}\) and \(\mathfrak M_\infty\)
      are indiscernible by means of finite sets as models for
      \(\mathcal T_2\), and
      \(\mathfrak M_{\omega^{\omega^{\omega^\omega}}}\) and
      \(\mathfrak M_\infty\) as models for \(\mathcal T_3\).
\end{enumerate}
\end{hbtheorem}

\begin{hbdefinition}{def:p03-w03-p03-038-e61-definable}{Definable ordinals}
An ordinal \(a\) is \emph{definable in a model} \(\mathfrak M\) for
\(\mathcal T_i\) (or \(T_i\)) iff some formula \(\Psi\in\mathcal T_i\)
(\(\in T_i\)) with one free variable satisfies
\(\operatorname{stsf}_{\mathfrak M}\Psi(b)\equiv b=a\) for every element
\(b\) of the model \cite[p.~140]{Ehrenfeucht1961}. It is \emph{definable
in} \(\mathcal T_i\) (\(T_i\)) if it is definable in
\(\mathfrak M_\infty\). \(D(\mathfrak M,\mathcal T_i)\) and
\(D(\mathfrak M,T_i)\) are the sets of ordinals definable in
\(\mathfrak M\) for \(\mathcal T_i\) and for \(T_i\). Put
\(a_1\defeq\omega^\omega\), \(a_2\defeq\omega^{\omega^\omega}\),
\(a_3\defeq\omega^{\omega^{\omega^\omega}}\). The source records
Tarski's conjecture that \(x\) is definable in \(T_i\) iff \(x<a_i\),
\(i=1,2,3\) \cite[p.~140]{Ehrenfeucht1961}, and refers for Tarski's
results to his abstract in the Notices of the American Mathematical
Society (its reference~[6], p.~141).
\end{hbdefinition}

\begin{hbtheorem}{thm:p03-w03-p03-038-e61-below}{proved=annex:p03-w03-p03-038-e61-below}{Ordinals below \(a_i\) are definable}
For \(i=1,2,3\): (\(\bullet\)) every \(x<a_i\) is definable in \(T_i\);
and (\(\bullet\bullet\)) \(D(\mathfrak M,T_i)\subseteq
D(\mathfrak M,\mathcal T_i)\) for every model \(\mathfrak M\) for
\(\mathcal T_i\).
\end{hbtheorem}

The source calls (\(\bullet\)) not difficult and (\(\bullet\bullet\))
obvious, and proves neither \cite[p.~140]{Ehrenfeucht1961}. The annex
defines, in \(T_1\), the multiples of \(\omega^k\) by induction on \(k\)
and builds every \(x<\omega^\omega\) from \(0\) by steps
\(y\mapsto y+\omega^k\); for \(T_2\) and \(T_3\) it relativizes these
definitions to the ordinals \(\omega^\gamma\), which \(\sigma\) defines as
the nonzero ordinals closed under addition, reading \(\nu\) as \(\nu\) for
\(T_2\) and \(\sigma\) as \(\mu\) for \(T_3\).

\begin{hblemma}{thm:p03-w03-p03-038-e61-initial}{proved=annex:p03-w03-p03-038-e61-initial}{The initial-segment lemma}
Let \(i\in\{1,2,3\}\) and \(a\ge1\). If every \(x<a\) belongs to
\(D(\mathfrak M_\infty,\mathcal T_i)\), and \(\mathfrak M_a\) and
\(\mathfrak M_\infty\) are indiscernible by means of finite sets as
models for \(\mathcal T_i\), then
\(D(\mathfrak M_\infty,\mathcal T_i)=\{x:x<a\}\)
\cite[p.~140]{Ehrenfeucht1961}.
\end{hblemma}

The source proves the lemma by considering the least definable \(x\) that
is not below \(a\) and a formula defining it, and stops there. The annex
completes the argument: each formula that defines an ordinal in
\(\mathfrak M_\infty\) defines an ordinal of \(\mathfrak M_a\), the
resulting map is strictly increasing, and a strictly increasing map from
\(\{x:x\le a\}\) into \(\{x:x<a\}\) cannot exist.

\begin{hbtheorem}{thm:p03-w03-p03-038-e61-thm18}{proved=annex:p03-w03-p03-038-e61-thm18}{Theorem~18 of the source}
For \(i=1,2,3\) and every ordinal \(x\), the following are equivalent
\cite[p.~141]{Ehrenfeucht1961}: (a) \(x\) is definable in
\(\mathcal T_i\); (b) \(x\) is definable in \(T_i\); (c) \(x<a_i\). In
particular Tarski's conjecture of
\cref{def:p03-w03-p03-038-e61-definable} holds.
\end{hbtheorem}

\begin{hbexample}{def:p03-w03-p03-038-e61-worked}{A first move in \(H_2(\mathfrak M_{\omega^2},\mathfrak M_{\omega^2\cdot2})\)}
Here \(\omega^2\equiv\omega^2\cdot2\pmod{\omega^2}\) with \(d=1\),
\(d'=2\), \(e=0\), so II wins \(H_2\). Let I choose \(\omega+5\) and
\(\omega\cdot3\) in \(\mathfrak M_{\omega^2}\); the proof of Theorem~12
works with \(k=1\). The segment types are \(\omega+5=\omega\cdot1+5\);
then \(\omega\cdot2=\omega\cdot2+0\), since
\((\omega+5)+1+\omega\cdot2=\omega\cdot3\); and \(\omega^2\), since
\(\omega\cdot3+1+\omega^2=\omega^2\). II answers
\(b_1=\omega\cdot\operatorname{sgn}1+5=\omega+5\) and
\(b_2=b_1+1+\omega\cdot\operatorname{sgn}2=\omega\cdot2\). The answering
segments have types \(\omega+5\), \(\omega\) and \(\omega^2\cdot2\), and
each is congruent modulo \(\omega\) to its partner: equal in the first
case, both infinite with remainder \(0\) in the other two. By
Theorem~12 for \(n=1\), II wins \(H_1\) on each pair of segments, and
Theorem~11 assembles these into the remaining move of \(H_2\).

For definability: \(\omega\) is the least nonzero ordinal satisfying the
formula \(D_1\) of the annex, which says in \(T_1\) that an ordinal is
\(0\) or a limit.
Relativized to the ordinals closed under addition, the same formula
defines \(\omega^\omega\) in \(T_2\), the least such ordinal that is a
limit of such ordinals. By Theorem~18, \(\omega^\omega\) is not definable
in \(T_1\), not even with formulas about finite sets.
\end{hbexample}

\begin{hbwarning}{def:p03-w03-p03-038-e61-converse}{Congruence is sufficient, not necessary}
\(\omega\not\equiv\omega+1\pmod\omega\), yet II wins
\(H_1(\mathfrak M_\omega,\mathfrak M_{\omega+1})\); and I wins
\(H_2(\mathfrak M_\omega,\mathfrak M_{\omega+1})\)
(\cref{thm:p03-w03-p03-038-e61-converse}). So Theorem~12 has no converse,
and one more round separates the two structures. Theorem~11 is stated
and used for \(\mathcal T_1\) only: its proof splits a structure into
segments, which respects the order but not addition or multiplication;
the passage to \(\mathcal T_2\) and \(\mathcal T_3\) goes through the
exponent maps of Theorems~14 and~16 instead.
\end{hbwarning}

\begin{hblemma}{thm:p03-w03-p03-038-e61-converse}{proved=annex:p03-w03-p03-038-e61-converse}{The ordinals \(\omega\) and \(\omega+1\)}
For \(\mathcal T_1\): \(\omega\not\equiv\omega+1\pmod\omega\);
\(W_1(\mathfrak M_\omega,\mathfrak M_{\omega+1})\) holds; and
\(W_2(\mathfrak M_\omega,\mathfrak M_{\omega+1})\) fails.
\end{hblemma}

\emph{Bridge to this section.} Player~I is Spoiler and Player~II is
Duplicator. A model \(\mathfrak M_a\) for \(\mathcal T_1\), restricted to
its individuals, is the ordinal \(a\) as a linear order \((a,<)\); for
\(\mathcal T_2\) and \(\mathcal T_3\) it is \(a\) with the graphs of
addition and multiplication restricted to \(a\). The relation
\(W_j\) says that Duplicator wins the remaining \(j\) rounds of the game
in which Spoiler may choose finitely many elements per round; a round of
\(G_n\) is such a round with one element, so \(W_n(\emptyset)\) implies a
Duplicator win in \(G_n\) of
\cref{def:p03-w03-p03-038-positions}. Elementary indiscernibility of
\(\mathfrak M_a\) and \(\mathfrak M_b\) for \(T_1\) is elementary
equivalence \((a,<)\equiv(b,<)\), so Theorem~13 gives it whenever
\(a\equiv b\pmod{\omega^\omega}\); for instance
\((\omega^\omega+5,<)\equiv(\omega^\omega\cdot2+5,<)\).
``Indiscernible by means of finite sets'' is agreement on all sentences of
\(\mathfrak H(A)\), a stronger condition than elementary equivalence.
\emph{Source route and printed slips.} The source proves Theorem~12 by
induction on \(n\), choosing II's first answers from the Cantor division
of segment types by \(\omega^k\); the annex follows that route. It
measures segments as half-open intervals \(\{z:x\le z<y\}\), while
Theorem~11 uses open segments; the annex computes the types of open
segments, which leaves the construction unchanged. The first display for
the \(b_i\) prints a factor \(a\) before \(d_1\); the construction reads
\(\omega^k\operatorname{sgn}c_1+d_1\), as in the other displays. In the
statement of Theorem~14 the exponents are listed as
\(a_1,\ldots,a_n\) with coefficients \(n_1,\ldots,n_k\); Theorem~15 is
said to follow from Theorems~12 and~13, and Theorem~16 refers to
``Theorem~13'' for \(c,d\) and for its proof, where the content used is
Theorem~14. Theorem~18 is printed for ``\(i=1,2,2\)'', where the
statement concerns \(i=1,2,3\). The proof of Theorem~18 applies
Theorems~13, 15 and~17 to the class of all ordinals, which those theorems,
being about two ordinals, do not literally cover; the class forms in
\cref{thm:p03-w03-p03-038-e61-thm12,thm:p03-w03-p03-038-e61-thm13,thm:p03-w03-p03-038-e61-thm14,thm:p03-w03-p03-038-e61-thm15}
supply that step by the same arguments, with \(W_j\) in place of
strategies so that no choice over a proper class is needed. The source
cites a summary by Feferman for the segment argument of Theorem~11 (its
footnote on p.~138).

\subsection{Closed games, longer games, and inexpressibility}
\label{sec:p03-w03-p03-038-long-games}

The finite game of fixed length cannot see past finitely many rounds.
Three extensions lengthen it: a game of length
\(\omega\), a game whose length is a finite number chosen by Spoiler below
an ordinal clock, and a game of transfinite length. Throughout this
subsection the signature is relational and need not be finite.

\begin{hbdefinition}{def:p03-w03-p03-038-games}{Longer comparison games}
Let \(\mathcal A,\mathcal B\) be structures and \((\bar a,\bar b)\) a
position.
\begin{enumerate}
\item In \(G_\omega(\mathcal A,\bar a;\mathcal B,\bar b)\) rounds
      \(0,1,2,\ldots\) are played as in \cref{def:p03-ef-game}. Duplicator
      wins an infinite play iff every finite stage of it is a partial
      isomorphism.
\item For an ordinal \(\alpha\), the \emph{clock game}
      \(G^\alpha(\mathcal A,\bar a;\mathcal B,\bar b)\) starts with clock
      value \(\alpha\). If the clock is \(0\) the play ends. Otherwise
      Spoiler chooses a smaller ordinal \(\beta\) as the new clock value
      together with an element of either structure, and Duplicator answers
      in the other. Every play is finite because ordinals admit no infinite
      descending sequence. Duplicator wins iff the final position is a
      partial isomorphism.
\item For an ordinal \(\theta\), the \emph{long game}
      \(G_\theta(\mathcal A,\bar a;\mathcal B,\bar b)\) has one round for
      each \(\xi<\theta\). The position at stage \(\xi\) is the initial
      position together with the pairs of all rounds before \(\xi\); at a
      limit stage it is therefore the union of the earlier positions.
      Duplicator wins iff the position at every stage
      \(\xi\le\theta\) is a partial isomorphism.
\end{enumerate}
A position of any length, finite or transfinite, is a \emph{partial
isomorphism} when every atomic formula, evaluated at finitely many of its
coordinates on the two sides, receives the same truth value on both sides;
for finite positions this is the definition of
\cref{def:p03-w03-p03-038-positions}. Thus a position is a partial
isomorphism iff each of its finite subpositions is one.
\end{hbdefinition}

For finite \(\theta=k\), \(G_k\) is the pointed game of
\cref{def:p03-w03-p03-038-positions}, and \(G_\omega\) is the long game of
length \(\omega\). For finite \(\alpha=k\), Duplicator wins \(G^k\) iff she
wins \(G_k\) (\cref{thm:p03-w03-p03-038-clock-comparison}).

\begin{hblemma}{thm:p03-w03-p03-038-closed}{proved=annex:p03-w03-p03-038-closed}{The winning condition is closed}
\begin{enumerate}
\item In every long game \(G_\theta\) whose initial position is a partial
      isomorphism, if Duplicator loses a play then there is a least stage
      at which the position fails to be a partial isomorphism, and that
      stage is a successor ordinal.
\item Give the set of moves (pairs of elements) the discrete topology and
      the set of infinite plays of \(G_\omega\) the product topology. The
      set of plays won by Duplicator is closed; the set won by Spoiler is
      open.
\item \(G_\omega(\mathcal A,\bar a;\mathcal B,\bar b)\) is determined: if
      Spoiler has no winning strategy, Duplicator has one. The proof uses
      the Axiom of Choice to select strategies and answers.
\end{enumerate}
\end{hblemma}

Part~(3) is the closed-game case of the Gale--Stewart theorem
\cref{thm:p06-gale-stewart}. That theorem is stated for moves in \(\omega\);
here moves range over arbitrary sets, so the proof is written out in the
annex. That argument does not extend to \(G_\theta\) for limit
\(\theta>\omega\): it keeps every position of the play inside the set of
positions from which Spoiler has no winning strategy, and it gives no
information about whether the union of such positions at a limit stage
again lies in that set. No
determinacy claim is made here for long games in general. The long games of
\cref{thm:p03-w03-p03-038-phi-games} are determined because a winning
strategy is constructed for one player in each case.

\begin{hbtheorem}{thm:p03-w03-p03-038-omega-comparison}{proved=annex:p03-w03-p03-038-omega-comparison}{The game of length \(\omega\) and \(L_{\infty\omega}\)}
Duplicator has a winning strategy in
\(G_\omega(\mathcal A,\bar a;\mathcal B,\bar b)\) iff some back-and-forth
family contains \((\bar a,\bar b)\). In particular Duplicator wins
\(G_\omega(\mathcal A,\mathcal B)\) iff \(\mathcal A\) and \(\mathcal B\)
are potentially isomorphic, iff they satisfy the same
\(L_{\infty\omega}\)-sentences.
\end{hbtheorem}

\begin{hbtheorem}{thm:p03-w03-p03-038-clock-comparison}{proved=annex:p03-w03-p03-038-clock-comparison}{Clock games and ordinal approximants}
For every ordinal \(\alpha<\omega_1\), Duplicator has a winning strategy in
\(G^\alpha(\mathcal A,\bar a;\mathcal B,\bar b)\) iff
\((\mathcal A,\bar a)\equiv_\alpha(\mathcal B,\bar b)\) in the sense of
\cref{def:p03-morley-approximants}. In that case the tuples agree on every
infinitary formula of quantifier rank at most \(\alpha\). For finite
\(\alpha=k\), Duplicator wins \(G^k\) iff she wins \(G_k\). If Duplicator
wins \(G_\omega(\mathcal A,\bar a;\mathcal B,\bar b)\), she wins every
\(G^\alpha(\mathcal A,\bar a;\mathcal B,\bar b)\).
\end{hbtheorem}

\begin{hbtheorem}{thm:p03-w03-p03-038-long-comparison}{proved=annex:p03-w03-p03-038-long-comparison}{The game of length \(\omega_1\)}
\begin{enumerate}
\item If Duplicator wins \(G_{\omega_1}(\mathcal A,\mathcal B)\), she wins
      \(G_\omega(\mathcal A,\mathcal B)\); hence
      \(\mathcal A\) and \(\mathcal B\) satisfy the same
      \(L_{\infty\omega}\)-sentences.
\item If \(\mathcal A\cong\mathcal B\), Duplicator wins
      \(G_\theta(\mathcal A,\mathcal B)\) for every \(\theta\).
\item If \(|\dom\mathcal A|,|\dom\mathcal B|\le\aleph_1\) and Duplicator wins
      \(G_{\omega_1}(\mathcal A,\mathcal B)\), then
      \(\mathcal A\cong\mathcal B\).
\end{enumerate}
\end{hbtheorem}

For structures of size at most \(\aleph_1\), then, the game of length
\(\omega_1\) is the isomorphism test, while the game of length \(\omega\) is
the \(L_{\infty\omega}\) test. \Cref{thm:p03-w03-p03-038-phi-games} gives
structures on which the two tests differ.

\begin{hbexample}{def:p03-w03-p03-038-countability}{Countability and completeness}
All finite order-preserving maps between \((\mathbb Q,<)\) and
\((\mathbb R,<)\) form a back-and-forth family by
\cref{thm:p03-rn-dlo-extension}, so Duplicator wins
\(G_\omega((\mathbb Q,<),(\mathbb R,<))\) and the two orders satisfy the same
\(L_{\infty\omega}\)-sentences (\cref{thm:p03-w03-p03-038-omega-comparison}).
One order is countable and the other is not; the other is Dedekind complete
and the first is not. Hence no \(L_{\infty\omega}\)-sentence in the
signature \(\{<\}\) is true in exactly the countable linear orders, and none
is true in exactly the Dedekind complete ones. The game of length
\(\omega_1\) separates them: Spoiler lists \(\mathbb Q\) in the first
\(\omega\) rounds and then chooses an element of \(\mathbb R\) outside the
countable range built so far. Every answer in \(\mathbb Q\) has already
been paired with a different real, so the next position is not injective
and Duplicator loses.
\end{hbexample}

\subsection{Clubs and the set-theoretic conditions of game arguments}
\label{sec:p03-w03-p03-038-clubs}

The long game of length \(\omega_1\) is analysed through the stages at
which a play has exhausted an initial segment of both structures. Those
stages form a club, and a stationary set of ordinals meets every club. This
subsection fixes \(\kappa=\omega_1\) in \cref{def:p06-club-stationary}.

\begin{hbdefinition}{def:p03-w03-p03-038-lim}{Countable limit ordinals and closure points}
Let \(\mathrm{Lim}\) be the set of nonzero limit ordinals below
\(\omega_1\). For \(g:\omega_1\to\omega_1\), an ordinal
\(\delta\in\mathrm{Lim}\) is a \emph{closure point} of \(g\) when
\(g(\gamma)<\delta\) for every \(\gamma<\delta\).
\end{hbdefinition}

\begin{hblemma}{thm:p03-w03-p03-038-club-facts}{proved=annex:p03-w03-p03-038-club-facts}{Club facts used by the game proofs}
\begin{enumerate}
\item \(\omega_1\) is regular: every countable set of countable ordinals
      has a countable supremum. This uses countable choice.
\item \(\mathrm{Lim}\) is club in \(\omega_1\).
\item The intersection of two clubs is club; hence every club is
      stationary, and a stationary set meets every club.
\item For every \(g:\omega_1\to\omega_1\), the closure points of \(g\)
      form a club.
\item The set \(\mathrm{Lim}'\) of limits of nonzero limit ordinals,
      \(\{\delta\in\mathrm{Lim}:\sup(\mathrm{Lim}\cap\delta)=\delta\}\), is
      club; the set \(\mathrm{Lim}\setminus\mathrm{Lim}'\) is unbounded and
      nonstationary.
\end{enumerate}
\end{hblemma}

\begin{hbdefinition}{def:p03-w03-p03-038-phi}{The orders \(\Phi(S)\)}
For \(S\subseteq\mathrm{Lim}\) and \(\alpha<\omega_1\), let
\(Q_\alpha(S)\defeq\mathbb Q\) if \(\alpha\notin S\) and
\(Q_\alpha(S)\defeq\{q\in\mathbb Q:q\ge0\}\) if \(\alpha\in S\). Let
\(\Phi(S)\) be the set \(\{(\alpha,q):\alpha<\omega_1,\ q\in Q_\alpha(S)\}\)
ordered lexicographically: \((\alpha,q)<(\beta,q')\) iff \(\alpha<\beta\),
or \(\alpha=\beta\) and \(q<q'\). For \(\delta\le\omega_1\), write
\(\Phi(S)\mathbin{\upharpoonright}\delta\) for the initial segment
\(\{(\alpha,q)\in\Phi(S):\alpha<\delta\}\), and call
\(\{\alpha\}\times Q_\alpha(S)\) the \emph{block} \(\alpha\). Write
\(S\mathbin{\triangle}T\) for the symmetric difference.
\end{hbdefinition}

\begin{hblemma}{thm:p03-w03-p03-038-phi-basic}{proved=annex:p03-w03-p03-038-phi-basic}{Basic properties of \(\Phi(S)\)}
Let \(S\subseteq\mathrm{Lim}\).
\begin{enumerate}
\item \(\Phi(S)\) is a dense linear order without endpoints, of
      cardinality \(\aleph_1\).
\item For \(\delta<\omega_1\), \(\Phi(S)\mathbin{\upharpoonright}\delta\)
      is countable.
\item For \(\delta\in\mathrm{Lim}\), the final segment
      \(\Phi(S)\setminus\Phi(S)\mathbin{\upharpoonright}\delta\) has a least
      element iff \(\delta\in S\), and that element is then \((\delta,0)\).
\end{enumerate}
\end{hblemma}

\begin{hbtheorem}{thm:p03-w03-p03-038-phi-iso}{proved=annex:p03-w03-p03-038-phi-iso}{Isomorphism of \(\Phi(S)\) and \(\Phi(T)\)}
For \(S,T\subseteq\mathrm{Lim}\),
\(\Phi(S)\cong\Phi(T)\) iff \(S\mathbin{\triangle}T\) is nonstationary in
\(\omega_1\).
\end{hbtheorem}

\begin{hbtheorem}{thm:p03-w03-p03-038-phi-games}{proved=annex:p03-w03-p03-038-phi-games}{Games on \(\Phi(S)\) and \(\Phi(T)\)}
Let \(S,T\subseteq\mathrm{Lim}\).
\begin{enumerate}
\item Duplicator wins \(G_\omega(\Phi(S),\Phi(T))\); hence \(\Phi(S)\) and
      \(\Phi(T)\) satisfy the same \(L_{\infty\omega}\)-sentences.
\item If \(S\mathbin{\triangle}T\) is nonstationary, Duplicator wins
      \(G_{\omega_1}(\Phi(S),\Phi(T))\).
\item If \(S\mathbin{\triangle}T\) is stationary, Spoiler wins
      \(G_{\omega_1}(\Phi(S),\Phi(T))\): he lists both orders, and at a stage
      \(\delta\in S\mathbin{\triangle}T\) where the position is an
      isomorphism between the two initial segments of length \(\delta\),
      he chooses \((\delta,0)\) in the order whose final segment has a least
      element, and then a point below Duplicator's answer.
\end{enumerate}
\end{hbtheorem}

\begin{hbcorollary}{thm:p03-w03-p03-038-phi-inexpressible}{proved=annex:p03-w03-p03-038-phi-inexpressible}{Inexpressibility from a stationary set}
\(\Phi(\varnothing)\) and \(\Phi(\mathrm{Lim})\) satisfy the same
\(L_{\infty\omega}\)-sentences but are not isomorphic, and Spoiler wins the
game of length \(\omega_1\) between them. Hence no
\(L_{\infty\omega}\)-sentence in the signature \(\{<\}\) holds in exactly
the linear orders isomorphic to \(\Phi(\mathrm{Lim})\), even among dense
linear orders without endpoints of cardinality \(\aleph_1\).
\end{hbcorollary}

\begin{hbwarning}{def:p03-w03-p03-038-unbounded-not-enough}{Unbounded is not stationary}
The hypothesis in \cref{thm:p03-w03-p03-038-phi-games}(3) cannot be
weakened to unboundedness. By \cref{thm:p03-w03-p03-038-club-facts}(5),
\(S=\mathrm{Lim}\setminus\mathrm{Lim}'\) is unbounded but nonstationary, so
\(\Phi(S)\cong\Phi(\varnothing)\) by \cref{thm:p03-w03-p03-038-phi-iso},
although block \(\delta\) has a least element for unboundedly many
\(\delta\). Spoiler's attack at such a \(\delta\) succeeds only at stages
where the position has exhausted both initial segments of length
\(\delta\); those stages form a club, and the attack needs one of them in
\(S\mathbin{\triangle}T\).
\end{hbwarning}

\emph{Where the assumptions enter.} Regularity of \(\omega_1\), from
countable choice, makes the stages at which a play or an isomorphism has
exhausted both initial segments unbounded. Closure of the club is the
statement that such exhaustion passes to limits. Stationarity of
\(S\mathbin{\triangle}T\) is used exactly once, to find an exhausted stage
at which the two final segments differ. The game of length \(\omega\) never
reaches such a stage, which is why it cannot see the difference.

\emph{Relationships.} \Cref{thm:p03-w03-p03-038-omega-comparison} is the
game form of \cref{thm:p08-bvb501-potential-isomorphism}, and
\cref{thm:p03-w03-p03-038-clock-comparison} is the game form of the
approximants used for Morley's countable-model analysis in
\cref{sec:p03-morley-scott}. The finite-signature games of
\cref{thm:p03-w03-p03-038-pointed-ef} are the clock games with a finite
clock.
